\chapter{Appuis, charges et combinaisons}

\section{Conditions aux limites}

Les appuis sont définis par les blocages des six degrés de liberté d'un nœud :

\begin{itemize}
  \item translations : UX, UY, UZ ;
  \item rotations : RX, RY, RZ.
\end{itemize}

Le logiciel propose des types d'appuis prédéfinis, mais il est aussi possible
de personnaliser les blocages.

\section{Cas de charge}

Un cas de charge regroupe des actions appliquées au modèle. Les types usuels
sont :

\begin{itemize}
  \item charges permanentes ;
  \item charges d'exploitation ;
  \item vent ;
  \item neige ;
  \item séisme ;
  \item poids propre automatique.
\end{itemize}

\section{Charges nodales}

Une charge nodale s'applique directement sur un nœud. Elle peut contenir des
forces et des moments selon les six degrés de liberté.

\section{Charges réparties sur barres}

Une charge répartie s'applique sur une barre. Elle peut être saisie dans le
repère local ou global selon le cas. Les conventions de signe doivent être
vérifiées attentivement.

\section{Charges surfaciques}

Les charges surfaciques s'appliquent aux plaques et dalles. Elles sont ensuite
converties en efforts équivalents pour le solveur lorsque la formulation le
permet.

\section{Combinaisons}

Les combinaisons regroupent plusieurs cas de charge avec des coefficients. Le
logiciel peut générer certaines combinaisons selon EC0, ou permettre une saisie
manuelle.

\todoDoc{Ajouter un tableau d'exemple pour une combinaison G + Q + S.}
